\section{One-sided support queries from pooled collision bits}
\label{sec:one-sided}
The high-accuracy occupation query estimates a difference of two means.
For boundary location this is not always necessary. On a short line
outside an expanded component, the forward collision probability itself
has a known zero side. Randomizing the nominal center along that line
turns detection of a positive response into a one-sided support query.
The four-direction law is unchanged and its realized direction remains
unobserved. Only the spatial commands are redesigned.

Put $\alpha=\gamma+3/2$. In this section choose the fixed dyadic
compass step sufficiently small that
\begin{equation}\label{eq:line-step}
                  8t+\overline d_K<d_0,
\end{equation}
where $\overline d_K$ is a known upper bound for $\diam K$ strictly
less than $d_0$. This is a choice of control at a fixed prior scale,
not a shrinking flight length. It is available whenever the footprint
has the separation required for a stationary experiment. The earlier
results remain valid for their larger range of steps.

\subsection{Coarse support frames}
Write $P=C+(-K)$ and $p=p_P^{\rm lab}$. For a unit normal
$n=n(\theta)$ its support point is
$z_P(n)=p(\theta)n+p'(\theta)n^\perp$. The coarse acquisition and
interpolation already proved supply, at any prescribed sufficiently
small \emph{fixed} tolerance, a support approximation in $C^1$ and
hence approximate support points uniformly in $n$. To justify this
without presupposing the new query, apply Proposition~\ref{prop:jitter-query}
and Lemmas~\ref{lem:coarse}--\ref{lem:radial} at fixed numerical scales.
Their cost depends on the tolerance and the priors, but not on the
subsequent accuracy $e$. The construction concerns $P$ alone and uses
only bounds for $K$, not evaluations of its support function. It works
also for a single isolated reference body in a calibration scene.

The components $P$ are separated by at least
$g_*=d_0-\overline d_K>8t$. Choose small constants $b,w>0$, and a
coarse error smaller than both, as follows. All line segments below
must lie within distance $t+b+3w$ of their designated support point,
with $t+b+3w<g_*-t$. Require also $3w<t/(2\sqrt2)$.
If $r_P>0$ is a common interior rolling radius for the expanded
bodies, reduce $w$ so that $3w<r_P/4$ and
$\sqrt{3r_Pw}/2<b/4$. Obtain a coarse support approximation
$\widetilde p$ with
$\|\widetilde p-p\|_{C^1}<\min(w/4,b/4)$, reducing constants further
when needed. All these choices are at fixed prior scales.

Let $\mathcal D=\{te_1,-te_1,te_2,-te_2\}$. For each $n$ choose
one maximizer $a_*(n)$ of $-a\cdot n$ over $a\in\mathcal D$ and set
\begin{equation}\label{eq:compass-support}
 m(n)=-a_*\cdot n=t\max(|n_1|,|n_2|)\ge t/\sqrt2,
 \qquad s_P(n)=p(n)+m(n).
\end{equation}
A tie is resolved by a fixed ordering; no geometric sign information
is supplied by this choice. Set
$u_0(n)=\widetilde p'(\theta)-a_*\cdot n^\perp$.
At level $r$, draw a controller variable $U$ uniformly on $[-b,b]$
and use the nominal center
\begin{equation}\label{eq:line-command}
             X=r n+(u_0(n)+U)n^\perp.
\end{equation}
At this center execute the usual forward command with independent
$Z\sim j$ and $a\sim\rho$. Let $q_P(r,n)$ be the mean of its one
returned bit. The normal, line, and nominal center are controller
inputs; neither the actual start nor the random compass direction is
observed.

\begin{lemma}[A zero-side support response]\label{lem:line-response}
On the fixed bracket
$|r-(\widetilde p(n)+m(n))|\le w$, there are uniform constants $c,w_1>0$
such that
\begin{align}
 r>s_P(n)&\quad\Longrightarrow\quad q_P(r,n)=0,
                                               \label{eq:line-zero}\\
 0<d=s_P(n)-r<w_1&\quad\Longrightarrow\quad
                q_P(r,n)\ge c d^{\gamma+2}.     \label{eq:line-mass}
\end{align}
The bracket is chosen to lie in the stated collar. Other components
contribute no collisions to these commands. The assertions hold at
compass ties and uniformly over the unknown density class.
\end{lemma}
\begin{proof}
A collision with another component $D$ would require
$X\in P_D+[-a,0]$. But $X$ lies within $t+b+3w$ of the chosen
point of $P$, while $P_D$ is at distance at least $g_*$ from $P$.
The preceding constants make $\dist(X,P_D)>t$ for every possible
center on the line. Thus neither a start in $D$ nor a collision with
$D$ is possible.

For any $z\in K$, any $a\in\mathcal D$ and $0\le\lambda\le1$,
\[
 (X+z+\lambda a)\cdot n
 \ge r-p_{-K}(n)-m(n)
   =p_C^{\rm lab}(n)+r-s_P(n).
\]
If $r>s_P(n)$, the entire attempted segment lies beyond a supporting
line of $C$. This proves \eqref{eq:line-zero} pointwise. On the
interior side, $r=s_P(n)-d$ with $d<3w$, every start satisfies
\[
 (X+z)\cdot n\ge p_C^{\rm lab}(n)+m(n)-d
                         >p_C^{\rm lab}(n).
\]
Every start is therefore free. Conditional on the single compass
choice $a=a_*$, an endpoint $X+z+a_*\in C$ necessarily produces a
hit. Consequently
\begin{equation}\label{eq:line-positive-bound}
       q_P(r,n)\ge\frac14\E_U v_j(X+a_*).
\end{equation}
No independent direction-resolved mean is being measured: the factor
$1/4$ is a lower bound for one contribution to the same pooled bit.

Now $X+a_*$ has normal coordinate $p(n)-d$ and tangential coordinate
$\widetilde p'(\theta)+U$. The disk of radius $r_P$ centered at
$z_P(n)-r_Pn$ lies in $P$. When
\[
 |\widetilde p'(\theta)+U-p'(\theta)|
                            \le\tfrac12\sqrt{r_Pd},
\]
the distance of $X+a_*$ from the boundary of this disk is at least
$d/2$, for $d$ below the chosen fixed threshold. Indeed
$(r_P-d)^2+r_Pd/4\le(r_P-d/2)^2$ for $d\le r_P$.
The entire displayed interval of $U$ lies in $[-b,b]$ by the coarse
reserve and the choice of $w$. Its uniform probability is at least
$c_1\sqrt d$. Lemma~\ref{lem:cap-mass} gives
$v_j(X+a_*)\ge c_2d^\alpha$ on that interval. Substitution in
\eqref{eq:line-positive-bound} proves
$d^{\alpha+1/2}=d^{\gamma+2}$. This argument uses any maximizing
$a_*$ and therefore includes ties. Boundary conventions cause no
ambiguity in the strict zero side or the interior lens.
\end{proof}

\subsection{Detection, interpolation and the sample exponent}
\begin{proposition}[One-sided finite detection]\label{prop:line-query}
For each line bracket, $0<e<e_0$ and conditional failure allowance
$0<\varepsilon<1/4$, at most
\begin{equation}\label{eq:line-query-cost}
                 C e^{-(\gamma+2)}\log(C/\varepsilon)
\end{equation}
forward attempts give a label that is correct for $r>s_P(n)$ and
$r<s_P(n)-e$, with conditional probability at least $1-\varepsilon$.
A label is unrestricted in the remaining layer. Only the pooled bit
is used. Dyadic nominal centers with error $O(e)$ enlarge the layer
to $O(e)$ without changing this order.
\end{proposition}
\begin{proof}
Return one if at least one hit occurs in the prescribed batch, and
zero otherwise. A true zero response produces no hit. If
$q_P(r,n)\ge c e^{\gamma+2}$, the probability of no hit in $M$
independent attempts is at most
$(1-q_P)^M\le\exp(-cMe^{\gamma+2})$.
This proves \eqref{eq:line-query-cost} and the conditional assertion
with fresh preparations. In particular, an estimate of two means to
absolute error $e^{\gamma+3/2}$ is not required.

For an explicitly finite control law, replace $U$ by a uniform choice
on a power-of-two grid in $[-b,b]$ with spacing at most $c_3e$.
Every interval of length $c_1\sqrt d$ in the proof, for $d\ge e$,
still receives mass at least $c_4\sqrt d$. Round each nominal center
with Euclidean error at most $c_5e$. Its normal coordinate moves by
at most that amount, so the zero assertion holds beyond the enlarged
layer. On a shorter tangential interval the perturbed endpoint remains
in $P_{-c_6d}$, giving the same lower bound. This is a geometric
argument and asks for no upper or derivative bound on $j$. Certified
normal and support-frame evaluations of error $O(e)$ are absorbed in
these reserves. The fixed compass displacement itself remains exact.
\end{proof}

\begin{theorem}[An improved fixed-spread upper bound]
\label{thm:line-reconstruction}
For the stationary experiment on $\mathcal B_\eta$, with known
footprint $K$ and the step chosen as in \eqref{eq:line-step}, there is
a finite controller independent of $j$ giving the same geometric and
primitive discrete conclusions as Theorem~\ref{thm:stationary-jitter},
with probability at least $1-\delta$, and using at most
\begin{equation}\label{eq:line-total}
 C\nu^{-((\gamma+2)s+1)/(s-2)}
                   \log(C/\nu)\log(C/(\nu\delta))
\end{equation}
attempted bits. For $\gamma=0$ the power is $(2s+1)/(s-2)$.
The initial coarse phase uses reciprocal commands; the fine phase
uses forward pooled commands with randomized nominal centers on short
lines. The footprint remains fixed. Nominal precision is
$O(\nu^{s/(s-2)})$ and no additional arbitrary position, angle or
flight-time error is allowed in this stationary model.

There are $O(h^{-1}\log(C/h))$ line-test batches, where
$h\asymp\nu^{1/(s-2)}$. These are not that many distinct centers:
with $e\asymp h^s$ the finite tangential grids have at most
$C h^{-1}e^{-1}\log(C/h)$ possible centers, and at most one center
command is issued per attempted bit. The bound prices attempts,
not apparatus travel, footprint manufacture or numerical bit operations.
It is an improved upper bound, not a stationary-noise minimax theorem.
\end{theorem}
\begin{proof}
Condition on successful coarse frames. For $m\asymp h^{-1}$ equally
spaced normals, bisect the scalar support bracket with
Proposition~\ref{prop:line-query} at $e=c h^s$.
The relaxed bisection invariant of Lemma~\ref{lem:bisection} applies
to the threshold $s_P(n)$ exactly as it did to a radial value; no
monotonicity inside the layer is used. Subtract the \emph{known}
number $m(n)$ from each returned value. This supplies values of the
smooth support function $p_P$ with error $O(e)$. The function $m(n)$
need not be differentiable at compass ties: subtraction is done on
the sample values before interpolation, not after differentiating an
estimated nonsmooth envelope.

The degree-six interpolation and partition argument in
Lemma~\ref{lem:radial} applies directly to support functions. It gives
\[
 \|\widehat p_P-p_P\|_{C^j}
      \le C(h^{s-j}+eh^{-j}),\qquad 0\le j\le3.
\]
Subtract $p_{-K}$ and retain the positive curvature margin. The result
is an actual smooth convex output after the disclosed finite smoothing
in Theorem~\ref{thm:stationary-jitter}; the latter costs no new
collision bits. Its support error is $O(h^{s-2})$ in $C^2$ and
$O(h^s)$ in supremum norm. Apply Lemma~\ref{lem:locking} to these
original components. Exact rational relations and orbit counts are
then correct, while real generators and free area remain estimates.

The fixed coarse phase costs a prior-dependent multiple of
$\log(C/\delta)$. The fine phase uses at most
$Q_h=C h^{-1}\log(C/h)$ batches. Allocate each conditional failure
allowance $\delta/(2Q_h)$ and reserve $\delta/2$ for coarse recovery.
The union bound and \eqref{eq:line-query-cost} give
\[
 N\le C h^{-1}e^{-(\gamma+2)}
                           \log(C/h)\log(C/(h\delta)).
\]
Substitute $e=c h^s$ to obtain \eqref{eq:line-total}. Every launch,
including unsuccessful batches and solid starts in the coarse phase,
is charged. The nominal-grid and command counts follow from the
power-of-two tangential grid in each batch. Normal approximations with
error $O(e)$ change the sampled support values by $O(e)$ on the bounded
laboratory aperture, hence fit the same interpolation reserve. This
also supplies a finite dyadic description of each issued command.
\end{proof}

The gain is specific and does not imply that the old occupation
estimator has smaller variance. A zero-side response is detected at
cost inverse in its mean. Tangential randomization adds one half to
the lens exponent, yielding $\gamma+2$ instead of $2\gamma+3$ in
the fine-scale cost. The algorithm uses additional spatial command
locations and does not establish optimality. The two constructions
remain useful distinct inverses in the same pooled-bit information
category.

\subsection{A matching cost for a fixed support-line test}
\begin{proposition}[Physical depth exponent for one line]
\label{prop:line-depth-sharp}
There are disk components and one common stationary launch density
satisfying the stated footprint class for which the pooled line response
at inward depth $e$ has probability comparable to $e^{\gamma+2}$,
whereas at the support level it has probability zero. A procedure that
uses only repeated independent trials of this fixed line command and
distinguishes these two alternatives with error at most
$0<\delta<1/4$ in each world must use, under a deterministic sample cap,
at least
\[
                         c e^{-(\gamma+2)}\log(1/\delta)
\]
attempts. This is a lower bound for the stipulated line experiment,
not for controllers that may choose other lines or arbitrary commands.
\end{proposition}
\begin{proof}
Take a disk of radius $r>0$, a footprint disk of radius $R>0$, and
\[
 j(z)=c_\gamma(R^2-|z|^2)^\gamma\one_{|z|<R},\qquad
              c_\gamma=\frac{\gamma+1}{\pi R^{2\gamma+2}}.
\]
The density is normalized and satisfies the required boundary-distance
lower bound. Use $n=e_1$ and the line with first coordinate
$r+R+t-e$. For small $e<t$, every direction except $-te_1$ misses
because the entire segment has first coordinate greater than $r$.
Starts are free. For direction $-te_1$, the endpoint has first
coordinate at least $r-e>0$, so the segment hits precisely when its
endpoint is in the disk. Taking the tangential interval long enough
to include this shrinking contact lens gives
\[
 q(e)=\frac{c_\gamma}{8b}
 B\!\left(\frac12,\gamma+1\right)
 \int_0^e (2Rw-w^2)^{\gamma+1/2}
                 2\sqrt{2r(e-w)-(e-w)^2}\,\dd w.
\]
Here $B$ is the beta integral. Scaling $w=ev$ bounds the integral
above and below by positive multiples of $e^{\gamma+2}$; the
remaining beta integral is finite and strictly positive. At $e=0$
the probability is zero. Placing identical disks on a sufficiently
widely spaced lattice yields physical periodic examples with the
other components outside the protective command aperture. The two
alternatives can be realized by translating that table by $en$ while
holding the reference line and the common noise law fixed.

The lower bound remains valid when the tangential controller draws
are recorded. Couple those independent draws and all parameter-independent
test randomness in the two worlds. Under the zero alternative every bit
is zero. Conditional on the draws, let $w$ be the probability of the
all-zero word under the positive alternative, and let $\phi$ be the
probability that the test accepts the positive alternative on this word.
Then $0\le w\le1$ and $\E\phi\le\delta$, while
$\E w=(1-q(e))^M$ for $M$ independent repetitions of the same line
experiment. Its error under the positive alternative is at least
\[
 \E[(1-\phi)w]\ge (1-q(e))^M-\delta.
\]
Uniform error at most $\delta$ implies $(1-q(e))^M\le2\delta$.
For small $e$, $-\log(1-q(e))\le2q(e)\le C e^{\gamma+2}$,
which proves the bound for $0<\delta<1/4$.
A procedure with a deterministic cap and earlier stopping can be
padded with unused independent bits; the same argument still applies.
The lower bound deliberately excludes changing the command location.
\end{proof}
